\documentclass[11pt]{article}

\usepackage[margin=1in]{geometry}
\usepackage{amsmath,amssymb}
\usepackage{booktabs}
\usepackage{array}
\usepackage{graphicx}
\usepackage[hidelinks]{hyperref}
\graphicspath{{reports/}{./}}

\title{WeatherEngression Synthetic NARX Student-\(t\) Experiment}
\author{WeatherEngression Research Notes}
\date{\today}

\begin{document}
\maketitle

\section{Goal}

The experiment tests whether engression can recover the conditional
distribution \(P(Y \mid X=x)\) on a synthetic weather-like time-series problem
where the true conditional quantiles are known.  The setup compares two
engression variants:

\begin{enumerate}
  \item \textbf{Vanilla engression}: the upstream public-package engression
  model, trained with its standard Adam optimizer.
  \item \textbf{Classic-Adam regularized engression}: the same upstream model,
  but trained with \(\texttt{torch.optim.Adam}\) using weight decay
  \(\lambda=0.003\).  This is Adam weight decay, not decoupled AdamW.
\end{enumerate}

Both models were trained on CPU with the same architecture and data.  The
reported quantiles are empirical quantiles of \(800\) generated samples
\(\widehat{Y}=g_\theta(x,\varepsilon)\) for each held-out \(x\).

\section{Data Generation}

The synthetic data model is \(\texttt{narx\_student\_t}\).  The run generated
\(n=5000\) supervised examples.  Each example has one scalar target \(Y_t\) and
a lag-window input
\[
  X_t = (W_{t-L},\ldots,W_t),
\]
where \(L=12\).  The covariate vector has dimension \(d=12\), so the neural
network receives a flattened input of dimension
\[
  (L+1)d = 13\cdot 12 = 156.
\]

The first six coordinates of \(W_t\) are synthetic weather variables:
temperature, humidity, pressure, radiation, precipitation, and wind speed.
They are generated by a stable nonlinear autoregressive process with daily and
seasonal forcing.  The remaining six coordinates are additional autocorrelated
seasonal covariates.  These extra covariates affect \(Y_t\) through a known
positive scale factor \(s_{\mathrm{extra}}(X_t)\), so the full
\(d=12\)-dimensional input matters.

\subsection{Latent Weather Index}

The generator also saves a scalar summary \(\phi(X_t)\).  This is not given to
engression as a special input; it is a diagnostic and splitting variable
computed from the full lag window.  Let \(a\) be the six base-feature weights
\[
  a = (0.45,-0.35,0.20,0.50,-0.65,0.25),
\]
and let lag weights satisfy
\[
  \omega_\ell \propto \exp(-\mathrm{age}_\ell/8),
  \qquad \sum_{\ell=0}^{L} \omega_\ell = 1.
\]
Then
\[
  \phi(X_t)
  =
  \sum_{\ell=0}^{L}
  \omega_\ell\, a^\top W^{\mathrm{base}}_{t-L+\ell}.
\]

\subsection{Extra-Feature Scale}

The six extra covariates do not enter \(\phi(X)\).  They enter the target
through the positive multiplicative factor \(s_{\mathrm{extra}}(X)\).  Let
\[
  Z_t = X^{\mathrm{extra}}_t \in \mathbb{R}^{(L+1)\times 6}
\]
denote the extra-feature lag window.  The code summarizes \(Z_t\) using three
vectors:
\[
  z_{\mathrm{current}} = Z_{t,L},
  \qquad
  z_{\mathrm{recent}} = \frac{1}{L+1}\sum_{\ell=0}^{L} Z_{t,\ell},
  \qquad
  z_{\mathrm{tendency}} = Z_{t,L} - Z_{t,0}.
\]
For \(d=12\), there are six extra covariates.  The raw signed weight vector is
\[
  \widetilde{w}
  =
  (0.60,-0.72,0.84,-0.96,1.08,-1.20),
\]
and the actual weight vector is normalized:
\[
  w = \frac{\widetilde{w}}{\lVert \widetilde{w}\rVert_2}.
\]
Let \(w_{\mathrm{rev}}\) be the same vector in reverse coordinate order.  The
extra-feature score is
\[
  r_{\mathrm{extra}}(Z_t)
  =
  0.55\,z_{\mathrm{recent}}^\top w
  +0.30\,\tanh\!\left(z_{\mathrm{current}}^\top w_{\mathrm{rev}}\right)
  +0.15\,z_{\mathrm{tendency}}^\top w.
\]
The scale itself is
\[
  s_{\mathrm{extra}}(Z_t)
  =
  \exp\!\left(0.25\,\tanh(r_{\mathrm{extra}}(Z_t))\right).
\]
This guarantees \(s_{\mathrm{extra}}(Z_t)>0\).  In principle,
\[
  e^{-0.25}
  \le
  s_{\mathrm{extra}}(Z_t)
  \le
  e^{0.25},
\]
approximately \([0.779,1.284]\), although the realized range in this run was
smaller: \(0.903\) to \(1.108\), with mean \(0.995\).

The important consequence is that the extra covariates do not change the
diagnostic scalar \(\phi(X)\), but they do change the conditional law of \(Y\):
\[
  Y_t
  =
  s_{\mathrm{extra}}(Z_t)\,Y_t^{\mathrm{base}}.
\]
Since \(s_{\mathrm{extra}}(Z_t)\) is always positive, the reference conditional
quantiles transform by direct multiplication.

\subsection{Conditional Target Law}

The base target follows a heavy-tailed NARX conditional distribution:
\[
  Y_t^{\mathrm{base}} \mid X_t=x
  =
  \mu(x) + 0.8\,\sigma_G(x)\,T_5,
\]
where \(T_5\) is a standard Student-\(t\) random variable with \(5\) degrees of
freedom.  The conditional mean is
\[
\begin{aligned}
\mu(x)
={}&
0.55\sin(\phi(x))
+0.18\phi(x)^2
+0.35\tanh(R_{\mathrm{recent}})
-0.45\sqrt{\max(P_{\mathrm{recent}},0)}
\\
&+0.25H_{\mathrm{current}}\log(1+U_{\mathrm{current}})
-0.18\Delta B
+0.12\sin(T_{\mathrm{long}}),
\end{aligned}
\]
where \(R\) is radiation, \(P\) is precipitation, \(H\) is humidity,
\(U\) is wind speed, \(B\) is pressure, and \(T\) is temperature.  The
subscripts ``current'', ``recent'', ``long'', and the tendency
\(\Delta B\) are deterministic summaries of the lag window.

The Gaussian scale reused by the Student-\(t\) model is
\[
\sigma_G(x)
=
0.12
+0.50\,
\operatorname{sigmoid}\!\left(
-0.4
+0.9P_{\mathrm{recent}}
+0.8|\Delta B|
-0.5R_{\mathrm{current}}
+0.35U_{\mathrm{current}}
\right).
\]
The final target multiplies the base target by the positive extra-feature
scale:
\[
  Y_t = s_{\mathrm{extra}}(X_t)\,Y_t^{\mathrm{base}}.
\]
In this run \(s_{\mathrm{extra}}(X_t)\) ranged from \(0.903\) to \(1.108\), with
mean \(0.995\).

Because the target law is known, the true conditional quantiles are known:
\[
  Q_\alpha(Y_t\mid X_t=x)
  =
  s_{\mathrm{extra}}(x)
  \left[
  \mu(x) + 0.8\,\sigma_G(x)\,t_{5,\alpha}
  \right],
\]
where \(t_{5,\alpha}\) is the \(\alpha\)-quantile of a standard
Student-\(t_5\).  The experiments store \(Q_{0.05}\), \(Q_{0.50}\), and
\(Q_{0.95}\).

\section{Train/Test and Support Definitions}

\subsection{In-Support Diagnostic}

For the in-support diagnostic, \(4000\) rows are selected for training from a
random permutation of the \(5000\) rows.  Held-out rows are then selected from
the remaining rows only if their \(\phi(X)\) values lie inside the central
training \(\phi\)-range:
\[
  Q_{0.02}\{\phi(X_i): i\in \mathcal{I}_{\mathrm{train}}\}
  \le \phi(X_j) \le
  Q_{0.98}\{\phi(X_i): i\in \mathcal{I}_{\mathrm{train}}\}.
\]
The resulting in-support evaluation set has \(500\) rows.  In this run, the
in-support training range was approximately
\[
  [-1.329,\ 5.303],
\]
and the test range was
\[
  [-1.311,\ 5.295].
\]
Therefore no in-support test row was outside the scalar \(\phi(X)\) support
proxy.

\subsection{Out-of-Support Comparison}

For the OOS comparison, one shared training set was used for each model.  Rows
were sorted by \(\phi(X)\).  The central \(4000\) rows were used for training,
and the \(1000\) tail rows were used as candidate OOS rows:
\[
  \mathcal{I}_{\mathrm{train}}
  =
  \text{central }4000\text{ rows by }\phi(X).
\]
The training scalar support was approximately
\[
  \phi(X)\in[-0.040,\ 4.612].
\]
The candidate pool contained \(500\) lower-tail rows and \(500\) upper-tail rows.

Four OOS rules were then evaluated on the same candidate pool:

\begin{enumerate}
  \item \textbf{Scalar projection}: flag \(x\) if
  \(\phi(x)<\min_{\mathrm{train}}\phi(X)\) or
  \(\phi(x)>\max_{\mathrm{train}}\phi(X)\).
  \item \textbf{Marginal range}: flatten \(X\) to \(156\) features and flag
  \(x\) if any feature lies outside the training min/max range for that
  feature.
  \item \textbf{Marginal quantile}: flatten \(X\) and flag \(x\) if any feature
  lies outside the training interval
  \([Q_{0.01}, Q_{0.99}]\) for that feature.
  \item \textbf{kNN distance}: standardize flattened \(X\) using training
  column moments.  Then flag \(x\) if its nearest-neighbor distance to the
  training reference set exceeds the \(0.95\)-quantile of train-to-train
  nearest-neighbor distances.  The reference set size was \(2000\).
\end{enumerate}

When a rule produced more than \(500\) flagged rows, the evaluation subset was
capped at \(500\) rows by taking evenly spaced rows in chronological order.

\section{Model and Training Configuration}

Both models used:
\[
\text{layers}=3,\quad
\text{hidden dimension}=192,\quad
\text{noise dimension}=96,\quad
\eta=0.003,
\]
with batch size \(512\), \(120\) epochs, standardized inputs, and CPU training.
The regularized model additionally used Adam weight decay \(\lambda=0.003\).

\section{Reproducibility}

The reported runs are intended to be reproducible from the fixed seed
\[
  \text{seed}=2026.
\]
The code fixes NumPy and PyTorch seeds before fitting.  For the in-support
diagnostics, the train/test split uses seed \(2026+17=2043\), and prediction
sampling uses seed \(2026+101=2127\).  For the OOS comparison, the central
\(\phi(X)\) training split is deterministic after data generation; prediction
sampling again uses seed \(2127\), and the kNN reference subset uses seed
\(2026+211=2237\).

Exact bit-level reproducibility still assumes the same code, package versions,
and CPU numerical backend.  Small floating-point differences are possible if
PyTorch, BLAS, or the engression package version changes.

\section{Metrics}

For \(\alpha\in\{0.05,0.50,0.95\}\), let \(q_\alpha(x)\) be the true generator
quantile and \(\widehat{q}_\alpha(x)\) be the engression-predicted quantile.
The mean quantile MAE is
\[
  \frac{1}{3}\sum_{\alpha\in\{0.05,0.50,0.95\}}
  \frac{1}{m}\sum_{j=1}^{m}
  \left|\widehat{q}_\alpha(x_j)-q_\alpha(x_j)\right|.
\]
Predicted interval coverage is
\[
  \frac{1}{m}
  \sum_{j=1}^{m}
  \mathbf{1}\left\{
  Y_j\in[\widehat{q}_{0.05}(x_j),\widehat{q}_{0.95}(x_j)]
  \right\}.
\]
True interval coverage replaces the predicted quantiles with the generator
quantiles.  The width ratio is
\[
  \frac{
  \frac{1}{m}\sum_j
  \left(\widehat{q}_{0.95}(x_j)-\widehat{q}_{0.05}(x_j)\right)
  }{
  \frac{1}{m}\sum_j
  \left(q_{0.95}(x_j)-q_{0.05}(x_j)\right)
  }.
\]

\subsection{Result Table Columns}

The result tables use the following columns:

\begin{itemize}
  \item \textbf{Model}: the fitted engression variant.  ``Vanilla'' is the
  public-package baseline, while ``Adam weight decay'' is the same model
  trained with classic Adam weight decay \(\lambda=0.003\).
  \item \textbf{OOS rule}: the empirical out-of-support rule used to select
  held-out candidate rows.  This column appears only in OOS tables.
  \item \textbf{OOS \(n\)}: the number of candidate rows flagged by that OOS
  rule before applying the evaluation cap.
  \item \textbf{Eval \(n\)}: the number of flagged rows actually used for
  metric computation.  This is capped at \(500\) for large OOS subsets.
  \item \textbf{Mean qMAE}: the average absolute quantile error across
  \(q_{0.05}\), \(q_{0.50}\), and \(q_{0.95}\).
  \item \textbf{q05 MAE, q50 MAE, q95 MAE}: absolute error for the individual
  \(5\%\), \(50\%\), and \(95\%\) conditional quantiles.
  \item \textbf{Coverage}: empirical coverage of the model-predicted
  \(90\%\) interval
  \([\widehat{q}_{0.05}(X),\widehat{q}_{0.95}(X)]\).
  \item \textbf{True cov.}: empirical coverage of the true generator
  \(90\%\) interval \([q_{0.05}(X),q_{0.95}(X)]\) on the same realized
  \(Y\) values.  This is a finite-sample sanity check and should be near, but
  not exactly equal to, \(0.90\).
  \item \textbf{Width}: mean predicted \(90\%\) interval width,
  \(\widehat{q}_{0.95}(X)-\widehat{q}_{0.05}(X)\).
  \item \textbf{Width ratio}: mean predicted interval width divided by mean
  true generator interval width.  Values below \(1\) indicate intervals that
  are too narrow on average; values above \(1\) indicate intervals that are too
  wide on average.
\end{itemize}

\section{Results}

\subsection{In-Support Results}

\begin{center}
\small
\begin{tabular}{lrrrrrr}
\toprule
Model & Mean qMAE & q05 MAE & q50 MAE & q95 MAE & Coverage & Width ratio \\
\midrule
Vanilla & 0.260 & 0.301 & 0.200 & 0.278 & 0.698 & 0.681 \\
Adam weight decay & 0.157 & 0.157 & 0.126 & 0.190 & 0.866 & 1.007 \\
\bottomrule
\end{tabular}
\end{center}

\noindent\textbf{Table 1.} In-support diagnostic on \(500\) held-out rows.
True interval coverage on realized \(Y\) was \(0.920\) for both rows because
the same data split was used.
\vspace{0.8em}

The in-support posterior-band charts are placed after the numerical result
tables in Section~\ref{sec:diagnostic-figures}.  They are separated from the
tables so that the numerical comparison remains readable.

\subsection{OOS Results: Vanilla Engression}

\begin{center}
\small
\begin{tabular}{lrrrrrrr}
\toprule
OOS rule & OOS \(n\) & Eval \(n\) & Mean qMAE & Coverage & True cov. & Width & Width ratio \\
\midrule
Scalar projection & 1000 & 500 & 0.470 & 0.520 & 0.876 & 0.951 & 0.796 \\
Marginal range & 457 & 457 & 0.639 & 0.425 & 0.884 & 0.963 & 0.868 \\
Marginal quantile & 965 & 500 & 0.477 & 0.516 & 0.892 & 0.949 & 0.799 \\
kNN distance & 392 & 392 & 0.670 & 0.383 & 0.878 & 0.968 & 0.791 \\
\bottomrule
\end{tabular}
\end{center}

\noindent\textbf{Table 2.} OOS diagnostic for vanilla engression.  ``Width''
is the mean predicted \(90\%\) interval width.
\vspace{0.8em}

\subsection{OOS Results: Classic-Adam Regularized Engression}

\begin{center}
\small
\begin{tabular}{lrrrrrrr}
\toprule
OOS rule & OOS \(n\) & Eval \(n\) & Mean qMAE & Coverage & True cov. & Width & Width ratio \\
\midrule
Scalar projection & 1000 & 500 & 0.500 & 0.576 & 0.876 & 1.217 & 1.019 \\
Marginal range & 457 & 457 & 0.668 & 0.486 & 0.884 & 1.210 & 1.090 \\
Marginal quantile & 965 & 500 & 0.505 & 0.584 & 0.892 & 1.215 & 1.023 \\
kNN distance & 392 & 392 & 0.672 & 0.434 & 0.878 & 1.220 & 0.998 \\
\bottomrule
\end{tabular}
\end{center}

\noindent\textbf{Table 3.} OOS diagnostic for classic-Adam regularized
engression with \(\lambda=0.003\).
\vspace{0.8em}

\clearpage
\section{Diagnostic Figures}
\label{sec:diagnostic-figures}

The following charts show the in-support posterior-band diagnostics.  They are
not floated, so their order in the PDF follows their order in this section.
The blue band is the true generator \(90\%\) interval, and the orange band is
the engression-predicted \(90\%\) interval.

\subsection{Vanilla Engression}

\begin{center}
\includegraphics[width=0.92\textwidth]{in_support_vanilla_posterior_bands.png}
\end{center}

\noindent\textbf{Figure 1.} In-support posterior-band diagnostic for vanilla
engression.

\subsection{Classic-Adam Regularized Engression}

\begin{center}
\includegraphics[width=0.92\textwidth]{in_support_regularized_posterior_bands.png}
\end{center}

\noindent\textbf{Figure 2.} In-support posterior-band diagnostic for
classic-Adam regularized engression with \(\lambda=0.003\).

\clearpage
\section{Interpretation}

Regularized engression substantially improved interval width calibration.  In
support, the width ratio moved from \(0.681\) for vanilla to \(1.007\) for
classic-Adam regularized engression.  In OOS settings, the regularized model
also produced width ratios close to \(1\), while vanilla stayed near
\(0.79\)--\(0.87\).

However, matching average interval width did not fully solve calibration.  The
regularized model still under-covered realized \(Y\), especially under
\(\text{kNN}\)-defined OOS where predicted coverage was \(0.434\) versus true
coverage \(0.878\).  This suggests that the regularized model spreads its
posterior samples more appropriately on average, but still misplaces parts of
the conditional distribution under harder high-dimensional support shifts.

\section{Source Artifacts}

The code and saved metrics for this note are in:

\begin{itemize}
  \item \texttt{experiments/engression\_diagnostic.py}
  \item \texttt{experiments/oos\_comparison.py}
  \item \texttt{runs/engression\_diagnostics/narx\_student\_t/in-support/}
  \item \texttt{runs/oos\_comparisons/narx\_student\_t\_vanilla\_d12/}
  \item \texttt{runs/oos\_comparisons/narx\_student\_t\_regularized\_d12/}
\end{itemize}

\end{document}
